\begin{frame}{그림으로 읽기: 마진 경계와 서포트 벡터}
  \begin{columns}[T]
    \begin{column}{0.52\textwidth}
      \centering
      \includegraphics[width=\linewidth]{ch05_1_svm_margin.png}
    \end{column}
    \begin{column}{0.44\textwidth}
      \small
      \begin{itemize}
        \item \kb{마진 경계선}(f=$\pm1$)은 결정 경계(f=0)와 \textbf{평행}.
        \item 두 마진 경계선 \textbf{사이에 데이터가 없음}.
        \item 오렌지 원: \kb{서포트 벡터}만 마진 경계 위에.
      \end{itemize}
      \vskip0.5em
      {\footnotesize 코드의 마지막 두 줄이 경계에서 서포트 벡터까지의
      실제 거리 $= 1/\|w\|$임을 직접 재서, 마진 공식 $2/\|w\|$가
      추상적 수식이 아니라 \textbf{실제 거리}임을 확인.}
    \end{column}
  \end{columns}
\end{frame}
